\documentclass[10pt,a4paper]{amsart}
\usepackage{amsmath,amssymb,amsthm,graphicx}
\usepackage{verbatim}
\newenvironment{reference}{\comment}{\endcomment}
%\newenvironment{reference}{}{}
\newenvironment{slogan}{\comment}{\endcomment}
\newenvironment{history}{\comment}{\endcomment}

% For commutative diagrams we use Xy-pic
\usepackage[all]{xy}

% We use 2cell for 2-commutative diagrams.
\xyoption{2cell}
% Theorem environments.
%
\theoremstyle{plain}
\newtheorem{theorem}[subsection]{Theorem}
\newtheorem{proposition}[subsection]{Proposition}
\newtheorem{lemma}[subsection]{Lemma}

\theoremstyle{definition}
\newtheorem{definition}[subsection]{Definition}
\newtheorem{example}[subsection]{Example}
\newtheorem{exercise}[subsection]{Exercise}
\newtheorem{situation}[subsection]{Situation}

\theoremstyle{remark}
\newtheorem{remark}[subsection]{Remark}
\newtheorem{remarks}[subsection]{Remarks}

\numberwithin{equation}{subsection}

% Macros
%
\def\lim{\mathop{\mathrm{lim}}\nolimits}
\def\colim{\mathop{\mathrm{colim}}\nolimits}
\def\Spec{\mathop{\mathrm{Spec}}}
\def\Hom{\mathop{\mathrm{Hom}}\nolimits}
\def\Ext{\mathop{\mathrm{Ext}}\nolimits}
\def\SheafHom{\mathop{\mathcal{H}\!\mathit{om}}\nolimits}
\def\SheafExt{\mathop{\mathcal{E}\!\mathit{xt}}\nolimits}
\def\Sch{\mathit{Sch}}
\def\Mor{\mathop{\mathrm{Mor}}\nolimits}
\def\Ob{\mathop{\mathrm{Ob}}\nolimits}
\def\Sh{\mathop{\mathit{Sh}}\nolimits}
\def\NL{\mathop{N\!L}\nolimits}
\def\CH{\mathop{\mathrm{CH}}\nolimits}
\def\proetale{{pro\text{-}\acute{e}tale}}
\def\etale{{\acute{e}tale}}
\def\QCoh{\mathit{QCoh}}
\def\Ker{\mathop{\mathrm{Ker}}}
\def\Im{\mathop{\mathrm{Im}}}
\def\Coker{\mathop{\mathrm{Coker}}}
\def\Coim{\mathop{\mathrm{Coim}}}

% Boxtimes
%
\DeclareMathSymbol{\boxtimes}{\mathbin}{AMSa}{"02}

%
% Macros for moduli stacks/spaces
%
\def\QCohstack{\mathcal{QC}\!\mathit{oh}}
\def\Cohstack{\mathcal{C}\!\mathit{oh}}
\def\Spacesstack{\mathcal{S}\!\mathit{paces}}
\def\Quotfunctor{\mathrm{Quot}}
\def\Hilbfunctor{\mathrm{Hilb}}
\def\Curvesstack{\mathcal{C}\!\mathit{urves}}
\def\Polarizedstack{\mathcal{P}\!\mathit{olarized}}
\def\Complexesstack{\mathcal{C}\!\mathit{omplexes}}
% \Pic is the operator that assigns to X its picard group, usage \Pic(X)
% \Picardstack_{X/B} denotes the Picard stack of X over B
% \Picardfunctor_{X/B} denotes the Picard functor of X over B
\def\Pic{\mathop{\mathrm{Pic}}\nolimits}
\def\Picardstack{\mathcal{P}\!\mathit{ic}}
\def\Picardfunctor{\mathrm{Pic}}
\def\Deformationcategory{\mathcal{D}\!\mathit{ef}}
\usepackage{cleveref}
\crefname{originallemma}{original lemma}{original lemmas}
\Crefname{originallemma}{Original lemma}{Original lemmas}
\makeatletter
\newlabel{algebra-lemma-normal-domain-intersection-localizations-height-1}{{031T}{source}{}{sourcepointer-1}{}}
\newlabel{algebra-lemma-normal-domain-intersection-localizations-height-1@cref}{{[originallemma][0][]031T}{[1][0][]source}{}{}{}}
\makeatother
\crefname{originallemma}{original lemma}{original lemmas}
\Crefname{originallemma}{Original lemma}{Original lemmas}
\makeatletter
\newlabel{decent-spaces-lemma-henselian-local-ring-strict}{{0EMZ}{source}{}{sourcepointer-2}{}}
\newlabel{decent-spaces-lemma-henselian-local-ring-strict@cref}{{[originallemma][0][]0EMZ}{[1][0][]source}{}{}{}}
\makeatother
\crefname{originalequation}{original equation}{original equations}
\Crefname{originalequation}{Original equation}{Original equations}
\makeatletter
\newlabel{equation-c1}{{0BE7}{source}{}{sourcepointer-3}{}}
\newlabel{equation-c1@cref}{{[originalequation][0][]0BE7}{[1][0][]source}{}{}{}}
\makeatother
\crefname{originalreference}{original reference}{original references}
\Crefname{originalreference}{Original reference}{Original references}
\makeatletter
\newlabel{lemma-decent-irreducible-closed}{{0ENF}{source}{}{sourcepointer-4}{}}
\newlabel{lemma-decent-irreducible-closed@cref}{{[originalreference][0][]0ENF}{[1][0][]source}{}{}{}}
\makeatother
\crefname{originalreference}{original reference}{original references}
\Crefname{originalreference}{Original reference}{Original references}
\makeatletter
\newlabel{lemma-divisor-meromorphic-well-defined}{{02SH}{source}{}{sourcepointer-5}{}}
\newlabel{lemma-divisor-meromorphic-well-defined@cref}{{[originalreference][0][]02SH}{[1][0][]source}{}{}{}}
\makeatother
\crefname{originallemma}{original lemma}{original lemmas}
\Crefname{originallemma}{Original lemma}{Original lemmas}
\makeatletter
\newlabel{properties-lemma-normal-locally-finite-nr-irreducibles}{{0357}{source}{}{sourcepointer-6}{}}
\newlabel{properties-lemma-normal-locally-finite-nr-irreducibles@cref}{{[originallemma][0][]0357}{[1][0][]source}{}{}{}}
\makeatother
\crefname{originallemma}{original lemma}{original lemmas}
\Crefname{originallemma}{Original lemma}{Original lemmas}
\makeatletter
\newlabel{spaces-divisors-lemma-meromorphic-quasi-coherent}{{0EN6}{source}{}{sourcepointer-7}{}}
\newlabel{spaces-divisors-lemma-meromorphic-quasi-coherent@cref}{{[originallemma][0][]0EN6}{[1][0][]source}{}{}{}}
\makeatother
\crefname{originallemma}{original lemma}{original lemmas}
\Crefname{originallemma}{Original lemma}{Original lemmas}
\makeatletter
\newlabel{spaces-properties-lemma-describe-etale-local-ring}{{04KF}{source}{}{sourcepointer-8}{}}
\newlabel{spaces-properties-lemma-describe-etale-local-ring@cref}{{[originallemma][0][]04KF}{[1][0][]source}{}{}{}}
\makeatother
\title{The map from the Picard group to the divisor class group is injective for a normal integral algebraic space}
\author{The Stacks Project Authors}
\date{stacks-project:spaces-over-fields.tex:local-snapshot-20261012; GNU FDL 1.2 or later}
\begin{document}
\maketitle
\noindent\textit{Editorial note.} All delivered mathematical text is the author's own, copied from the pinned Stacks Project chapter file spaces-over-fields.tex; no statement or proof step was rewritten, reordered or removed.\par\medskip
\noindent\textit{Editorial note.} The counted proof is the author's complete argument: the reduction to a meromorphic section with zero divisor, the local computation over an integral affine chart with the identification of the relevant strict henselizations as discrete valuation rings, the comparison of local equations forcing the rational function to be a unit at every height one prime, and the closing use of the intersection theorem for normal domains for the section and its inverse.\par\medskip
\noindent\textit{Editorial note.} The label lemma-normal-c1-injective names the counted statement, and the proof environment immediately following it in the source is delivered verbatim as its proof.\par\medskip
\noindent\textit{Editorial note.} The counted statement is carried by the author's own lemma environment in the chapter source.\par\medskip
\noindent\textit{Editorial note.} Pre-open gate: python operations/agent-authoring/tools/precard\_gate.py --source spaces-over-fields.tex --label lemma-normal-c1-injective --start 1262 --end 1329 --local-id 28055 returned RESULT=OK (L28055 free, occupancy FREE, public interval clear). Candidate taken from sources/stacks-spaces-over-fields/candidates-halfA.json (fifth pool, pre-assigned by the Lead on 2026-10-13).\par\medskip
\noindent\textit{Editorial note.} The proof contains the sentence 'Our proof of the claim is a bit awkward as we don't yet have a lot of theory at our disposal; we suggest the reader skip the proof.' — a reading suggestion, while the proof environment then carries out the argument in full, so the statement is counted and the sentence is recorded here verbatim.\par\medskip
\noindent\textit{Editorial note.} The statement and proof are copied from the pinned chapter file; the cited inputs are kept as external pointers with their canonical Stacks tags.\par\medskip
\noindent\textit{Allowed standard prerequisites (editorial).}\begin{itemize}
\item Well definedness of the divisor of a meromorphic section. Same chapter, the lemma showing that the divisor associated to a meromorphic section is independent of the choices made. Supplies the manipulation of the divisor of the section t obtained by dividing by a rational function.
\item Normal local rings at codimension one points and their henselizations. Algebra, the lemma that a normal Noetherian domain is the intersection of its localizations at height one primes, together with Same chapter, the lemma identifying the strict henselization of the local ring at a codimension one point of a normal space as a discrete valuation ring. Supplies the discrete valuation ring computations and the conclusion that the rational function is a unit in the local ring and lies in the ring of the affine chart.
\item Meromorphic sections over an affine chart and trivializations. Divisors on Spaces, the lemma describing meromorphic sections over an affine open through the fraction field, together with the local triviality of invertible modules. Supplies the writing t restricted to the chart as f times a trivialization with f in the fraction field.
\item Local structure of a normal algebraic space. Properties of Spaces, the lemma that a normal algebraic space is locally a finite disjoint union of integral schemes, and Decent Spaces, the lemma describing henselian local rings of decent spaces. Supplies the reduction to an integral affine chart and the identification of the strict henselizations appearing in the computation.
\end{itemize}\medskip
\ifdefined\hypertarget\hypertarget{sourcepointer-1}{}\fi
\noindent\textit{Source reference.} Original lemma 031T: sources/stacks-chow/tags.txt, source lines 3996--3996. Prerequisite pointer invoked at proof line(s) 1325; it names a standard fact of a sibling chapter that the counted argument uses, and every step of the counted proof built on it is retained verbatim in the delivered body.\par\medskip
\ifdefined\hypertarget\hypertarget{sourcepointer-2}{}\fi
\noindent\textit{Source reference.} Original lemma 0EMZ: sources/stacks-chow/tags.txt, source lines 18324--18324. Prerequisite pointer invoked at proof line(s) 1305; it names a standard fact of a sibling chapter that the counted argument uses, and every step of the counted proof built on it is retained verbatim in the delivered body.\par\medskip
\ifdefined\hypertarget\hypertarget{sourcepointer-3}{}\fi
\noindent\textit{Source reference.} Original equation 0BE7: sources/stacks-spaces-over-fields/pinned-original-source.tex, source lines 1257--1257. Prerequisite pointer invoked at proof line(s) 1265; it names a standard fact of the same chapter that the counted argument uses, and every step of the counted proof built on it is retained verbatim in the delivered body.\par\medskip
\ifdefined\hypertarget\hypertarget{sourcepointer-4}{}\fi
\noindent\textit{Source reference.} Original reference 0ENF: sources/stacks-spaces-over-fields/pinned-original-source.tex, source lines 566--566. Prerequisite pointer invoked at proof line(s) 1317; it names a standard fact of the same chapter that the counted argument uses, and every step of the counted proof built on it is retained verbatim in the delivered body.\par\medskip
\ifdefined\hypertarget\hypertarget{sourcepointer-5}{}\fi
\noindent\textit{Source reference.} Original reference 02SH: sources/stacks-spaces-over-fields/pinned-original-source.tex, source lines 1168--1168. Prerequisite pointer invoked at proof line(s) 1277; it names a standard fact of the same chapter that the counted argument uses, and every step of the counted proof built on it is retained verbatim in the delivered body.\par\medskip
\ifdefined\hypertarget\hypertarget{sourcepointer-6}{}\fi
\noindent\textit{Source reference.} Original lemma 0357: sources/stacks-chow/tags.txt, source lines 4118--4118. Prerequisite pointer invoked at proof line(s) 1288; it names a standard fact of a sibling chapter that the counted argument uses, and every step of the counted proof built on it is retained verbatim in the delivered body.\par\medskip
\ifdefined\hypertarget\hypertarget{sourcepointer-7}{}\fi
\noindent\textit{Source reference.} Original lemma 0EN6: sources/stacks-chow/tags.txt, source lines 18331--18331. Prerequisite pointer invoked at proof line(s) 1292; it names a standard fact of a sibling chapter that the counted argument uses, and every step of the counted proof built on it is retained verbatim in the delivered body.\par\medskip
\ifdefined\hypertarget\hypertarget{sourcepointer-8}{}\fi
\noindent\textit{Source reference.} Original lemma 04KF: sources/stacks-chow/tags.txt, source lines 5971--5971. Prerequisite pointer invoked at proof line(s) 1308; it names a standard fact of a sibling chapter that the counted argument uses, and every step of the counted proof built on it is retained verbatim in the delivered body.\par\medskip
\begin{lemma}
\label{lemma-normal-c1-injective}
Let $S$ be a scheme. Let $X$ be a locally Noetherian integral algebraic space
over $S$. If $X$ is normal, then the map (\ref{equation-c1})
$\Pic(X) \to \text{Cl}(X)$ is injective.
\end{lemma}

\begin{proof}
Let $\mathcal{L}$ be an invertible $\mathcal{O}_X$-module whose
associated Weil divisor class is trivial. Let $s$ be a regular
meromorphic section of $\mathcal{L}$. The assumption means that
$\text{div}_\mathcal{L}(s) = \text{div}(f)$ for some
$f \in R(X)^*$. Then we see that $t = f^{-1}s$ is a regular
meromorphic section of $\mathcal{L}$ with
$\text{div}_\mathcal{L}(t) = 0$, see
Lemma \ref{lemma-divisor-meromorphic-well-defined}.
We claim that $t$ defines a trivialization of $\mathcal{L}$.
The claim finishes the proof of the lemma.
Our proof of the claim is a bit awkward as we
don't yet have a lot of theory at our dispposal; we suggest
the reader skip the proof.

\medskip\noindent
We may check our claim \'etale locally. Let $U \in X_\etale$ be affine
such that $\mathcal{L}|_U$ is trivial. Say $s_U \in \Gamma(U, \mathcal{L}|_U)$
is a trivialization. By
Properties, Lemma \ref{properties-lemma-normal-locally-finite-nr-irreducibles}
we may also assume $U$ is integral. Write $U = \Spec(A)$ as the spectrum of a
normal Noetherian domain $A$ with fraction field $K$.
We may write $t|_U = f s_U$ for some element $f$ of $K$, see
Divisors on Spaces, Lemma \ref{spaces-divisors-lemma-meromorphic-quasi-coherent}
for example. Let $\mathfrak p \subset A$ be a height one prime
corresponding to a codimension $1$ point $u \in U$ which maps
to a codimension $1$ point $\xi \in |X|$. Choose a trivialization
$s_\xi$ of $\mathcal{L}_\xi$ as in the beginning of this section.
Choose a geometric point $\overline{u}$ of $U$ lying over $u$.
Then
$$
(\mathcal{O}_{X, \xi}^h)^{sh} =
\mathcal{O}_{X, \overline{u}} =
\mathcal{O}_{U, u}^{sh} = (A_\mathfrak p)^{sh}
$$
see Decent Spaces, Lemmas
\ref{decent-spaces-lemma-henselian-local-ring-strict}
and
Properties of Spaces, Lemma
\ref{spaces-properties-lemma-describe-etale-local-ring}.
The normality of $X$ shows that all of these are
discrete valuation rings. The trivializations $s_U$ and $s_\xi$
differ by a unit as sections of $\mathcal{L}$ pulled back to
$\Spec(\mathcal{O}_{X, \overline{u}})$.
Write $t = f_\xi s_\xi$ with $f_\xi \in Q(\mathcal{O}_{X, \xi}^h)$.
We conclude that $f_\xi$ and $f$ differ by a unit
in $Q(\mathcal{O}_{X, \overline{u}})$.
If $Z \subset X$ denotes the prime divisor corresponding to $\xi$
(Lemma \ref{lemma-decent-irreducible-closed}), then
$0 = \text{ord}_{Z, \mathcal{L}}(t) =
\text{ord}_{\mathcal{O}_{X, \xi}^h}(f_\xi)$
and since $\mathcal{O}_{X, \xi}^h$ is a discrete valuation ring
we see that $f_\xi$ is a unit.
Thus $f$ is a unit in $\mathcal{O}_{X, \overline{u}}$
and hence in particular $f \in A_\mathfrak p$.
This implies $f \in A$ by Algebra, Lemma
\ref{algebra-lemma-normal-domain-intersection-localizations-height-1}.
We conclude that $t \in \Gamma(X, \mathcal{L})$.
Repeating the argument with $t^{-1}$ viewed as a meromorphic
section of $\mathcal{L}^{\otimes -1}$ finishes the proof.
\end{proof}

\end{document}
